\chapter{Bundle Maintenance and Curriculum Release}

\section{Purpose}
The offline student bundle is a reproducible distribution option, not a
different edition of uafR. It packages one source archive, the current student
manual, training scripts, integrity metadata, installation helpers, and an
acceptance test. Rebuilding the bundle for each approved package version keeps
student machines aligned and makes classroom support auditable.

\section{Release Inputs}
\begin{checklistbox}{Before Building}
\begin{itemize}
  \item Start from a reviewed repository commit with a clean working tree.
  \item Confirm that package tests and the release check pass without live web
        services.
  \item Build both public manuals and review their title pages, page counts,
        references, and institutional-mark placement.
  \item Run the training publication audit and offline training smoke tests.
  \item Confirm that no restricted data, credentials, unpublished project
        paths, or local usernames are present.
  \item Record the package version and source commit intended for the class.
\end{itemize}
\end{checklistbox}

\section{Build the Public Manuals}
From the repository root:

\begin{verbatim}
Rscript tools/build_training_manuals.R
Rscript tools/check_training_publication.R
\end{verbatim}

The build must produce:

\begin{itemize}
  \item \path{training/uafR_training_manual.pdf}
  \item \path{training/uafR_instructor_guide.pdf}
\end{itemize}

The student bundle includes only the student manual. The instructor guide is a
separate public companion and should not be placed in the student bundle as a
hidden answer key.

\section{Build the Student Bundle}
From the repository root:

\begin{verbatim}
Rscript tools/build_student_bundle.R
\end{verbatim}

An optional first argument chooses a different output directory:

\begin{verbatim}
Rscript tools/build_student_bundle.R /path/to/output
\end{verbatim}

The builder regenerates the student manual, builds the package source archive,
copies student-facing files, writes \path{MANIFEST.txt} and
\path{CHECKSUMS.csv}, and creates a versioned zip archive.

\section{Validate the Generated Bundle}
Unzip a fresh copy rather than testing only the source folder. Then run:

\begin{verbatim}
Rscript verify_bundle_integrity.R
Rscript preflight_check.R
Rscript install_uafR_from_bundle.R
Rscript run_student_acceptance_test.R
\end{verbatim}

For a strict release check, installation should occur in a clean temporary R
library on at least one supported operating system. The acceptance test remains
offline by default. A live database smoke test is optional and should never
determine whether a student bundle is structurally valid.

\section{Distribution Record}
Record the following before distribution:

\begin{itemize}
  \item Package version and source commit.
  \item Bundle zip filename, byte size, and checksum.
  \item Student-manual version, revision date, page count, and checksum.
  \item Date and platform used for the clean installation test.
  \item Known warnings or temporary service limitations.
  \item Approved distribution location and support contact.
\end{itemize}

\section{Updating an Existing Cohort}
Issue a new versioned bundle instead of replacing files inside a distributed
folder. Students should unzip the new bundle into a new directory and run
\path{update_uafR_from_bundle.R}. Keep the old bundle until active projects
have been rerun and the change is documented in each research log.

\begin{goldbox}{Release Boundary}
Public curriculum files may describe generic workflows and simulated examples.
Project-specific raw data, unpublished collaborator records, credentials,
internal schedules, and restricted storage locations must remain outside the
repository and outside the student bundle.
\end{goldbox}
